\thispagestyle{empty}

\thispagestyle{empty}
\begin{center}
\large{\thesisOrganization}
\end{center}

\vspace{0pt plus1fill}
\noindent%
\begin{tabularx}{\textwidth}{@{}lXr@{}}%
    & & \large{\textit{На правах рукописи}}\\
    % \IfFileExists{images/logo.pdf}{\includegraphics[height=2.5cm]{logo}}{\rule[0pt]{0pt}{2.5cm}}  & &
    \ifnumequal{\value{showperssign}}{0}{%
        \rule[0pt]{0pt}{1.5cm}
    }{
        \includegraphics[height=1.5cm]{personal-signature.png}
    }\\
\end{tabularx}

\vspace{0pt plus2fill} %число перед fill = кратность относительно некоторого расстояния fill, кусками которого заполнены пустые места
\begin{center}
\textbf {\Large \thesisAuthor}
\end{center}

\vspace{0pt plus1fill} %число перед fill = кратность относительно некоторого расстояния fill, кусками которого заполнены пустые места
\begin{center}
    \huge{\textbf{Правильные семейства функций \\ и порождаемые ими квазигруппы: \\ комбинаторные и алгебраические свойства}}

\vspace{0pt plus3fill} %число перед fill = кратность относительно некоторого расстояния fill, кусками которого заполнены пустые места
{\Large \textbf{\thesisSpecialtyNumber \; \thesisSpecialtyTitle}}

\vspace{0pt plus1.5fill} %число перед fill = кратность относительно некоторого расстояния fill, кусками которого заполнены пустые места
\Large{АВТОРЕФЕРАТ}\par
\large{диссертации на соискание учёной степени\par \thesisDegree}
\end{center}

\vspace{0pt plus4fill} %число перед fill = кратность относительно некоторого расстояния fill, кусками которого заполнены пустые места
{\centering\thesisCity~--- \thesisYear\par}

\newpage
% оборотная сторона обложки
\thispagestyle{empty}
\noindent Диссертация подготовлена на {\thesisInOrganization}.

\vspace{0.008\paperheight plus1fill}
\noindent%
\begin{tabularx}{\textwidth}{@{}lX@{}}
    \ifdefined\supervisorTwoFio
    Научные руководители:   & \textbf{\supervisorFio,}\par
                              \supervisorRegalia
                              \par
                              \vspace{0.013\paperheight}
                              \textbf{\supervisorTwoFio,}\par
                              \supervisorRegalia
                              \vspace{0.013\paperheight}\\
    \else
    Научный руководитель:   & \supervisorRegalia\par
                              \ifdefined\supervisorDead
                              \framebox{\textbf{\supervisorFio}}
                              \else
                              \textbf{\supervisorFio}
                              \fi
                              \vspace{0.013\paperheight}\\
    \fi
    Официальные оппоненты:  &
    \ifnumequal{\value{showopplead}}{0}{\vspace{10\onelineskip plus1fill}}{%
        \textbf{\opponentOneFio,}\par
        \opponentOneRegalia,\par
        \opponentOneJobPlace,
        \opponentOneJobPost

        \vspace{0.01\paperheight}
        \textbf{\opponentTwoFio,}\par
        \opponentTwoRegalia,\par
        \opponentTwoJobPlace, 
        \opponentTwoJobPost

    \ifdefined\opponentThreeFio
        \par
        \vspace{0.01\paperheight}
        \textbf{\opponentThreeFio,}\par
        \opponentThreeRegalia,\par
        \opponentThreeJobPlace, 
        \opponentThreeJobPost
    \fi
    }
    \vspace{0.013\paperheight} \\
    \ifdefined\leadingOrganizationTitle
    Ведущая организация:    &
    \ifnumequal{\value{showopplead}}{0}{\vspace{6\onelineskip plus1fill}}{%
        \leadingOrganizationTitle
    }%
    \fi
\end{tabularx}
%\vspace{0.008\paperheight plus1fill}

%\noindent 
Защита диссертации состоится \defenseDate~на~заседании диссертационного совета \defenseCouncilNumber~Московского государственного университета имени М. В. Ломоносова~по адресу: \defenseCouncilAddress.

%\vspace{0.008\paperheight plus1fill}
%\noindent 
E-mail: \defenseCouncilPhone.

%\vspace{0.008\paperheight plus1fill}
%\noindent 
С диссертацией можно ознакомиться в \synopsisLibrary

%\vspace{0.008\paperheight plus1fill}
%\noindent Отзывы на автореферат в двух экземплярах, заверенные печатью учреждения, просьба направлять по адресу: \defenseCouncilAddress, ученому секретарю диссертационного совета~\defenseCouncilNumber.

%\vspace{0.008\paperheight plus1fill}
%\noindent{
Автореферат разослан \synopsisDate.%}

\vspace{0.008\paperheight plus1fill}
\noindent%
\begin{tabularx}{\textwidth}{@{}%
>{\raggedright\arraybackslash}b{18em}@{}
>{\centering\arraybackslash}X
r
@{}}
    Ученый секретарь\par
    диссертационного совета \defenseCouncilNumber,\par
    \defenseSecretaryRegalia
    &
    \ifnumequal{\value{showsecrsign}}{0}{}{%
        \includegraphics[width=2cm]{secretary-signature.png}%
    }%
    &
    \defenseSecretaryFio
\end{tabularx}
